% Exemplo: tabela simples no padrão IBGE (booktabs + siunitx + \IBGEtab do abntex2)
% Compilar: latexmk -pdf tabela-simples.tex       Testado com: abntex2, TeX Live 2024 ou mais recente
% O trecho marcado com o id "tabelasimples" é o que as receitas e o livro exibem.
\documentclass[12pt, a4paper, brazilian]{abntex2}
\usepackage[T1]{fontenc}
%%BEGIN:tabelapreambulo
\usepackage{booktabs}                 % \toprule, \midrule, \bottomrule
\usepackage{siunitx}                  % colunas S: alinhamento decimal
% vírgula decimal e ponto nos milhares (o padrão do siunitx é o ponto decimal)
\sisetup{output-decimal-marker = {,}, group-separator = {.},
         group-minimum-digits = 4}
%%END:tabelapreambulo

\begin{document}

Os dados aparecem na \autoref{tab:simples}.

%%BEGIN:tabelasimples
\begin{table}[htbp]
\IBGEtab{%
  \caption{Usuários ativos e tempo médio de sessão por plataforma}%
  \label{tab:simples}%
}{%
  \begin{tabular}{l S[table-format=5.0]
                  S[table-format=2.1] S[table-format=2.1]}
    \toprule
    \textbf{Plataforma} & {\textbf{Usuários (un.)}}
      & {\textbf{Sessão (min)}} & {\textbf{Rejeição (\%)}} \\
    \midrule
    Web     & 12840 &  8,4 & 31,2 \\
    Android &  9317 & 11,9 & 24,7 \\
    iOS     &  4102 & 12,3 & 22,0 \\
    \bottomrule
  \end{tabular}%
}{%
  \fonte{Elaborado pelo próprio autor.}%
  \nota{Dados fictícios, apenas para exemplo.}%
}
\end{table}
%%END:tabelasimples

\end{document}
